\documentclass[11pt]{article}

\usepackage[spanish]{babel}
\usepackage[utf8]{inputenc}
\usepackage[T1]{fontenc}
\usepackage{geometry}
\usepackage{booktabs}
\usepackage{graphicx}
\usepackage{hyperref}
\usepackage{amsmath}
\usepackage{float}

\geometry{margin=2.5cm}

\title{Proyecto Integrador de Teoría del Riesgo\\
Dashboard de Riesgo Financiero para un Portafolio Internacional}
\author{Nombre del estudiante / grupo}
\date{\today}

\begin{document}

\maketitle

\begin{abstract}
Este trabajo desarrolla un tablero interactivo para el análisis de riesgo financiero de un portafolio compuesto por cinco acciones internacionales: Seven \& i Holdings, Alimentation Couche-Tard, FEMSA, BP y Carrefour. El proyecto integra descarga automática de datos desde APIs, análisis técnico, estudio de rendimientos, modelos ARCH/GARCH, estimación CAPM, medidas VaR/CVaR, optimización de Markowitz, señales automáticas y comparación contra benchmark junto con contexto macroeconómico.
\end{abstract}

\section{Introducción}
La gestión del riesgo financiero requiere integrar herramientas estadísticas, financieras y computacionales para apoyar la toma de decisiones. En este proyecto se construye un dashboard en Streamlit con actualización dinámica de datos y módulos especializados para evaluar diferentes dimensiones del riesgo.

\section{Datos y APIs}
Los precios históricos se obtuvieron mediante Yahoo Finance usando la librería \texttt{yfinance}. Para el contexto macroeconómico se utilizó la API de FRED, incluyendo la serie de tasa libre de riesgo a 3 meses, inflación y tipo de cambio COP/USD.

\section{Activos analizados}
\begin{itemize}
    \item Seven \& i Holdings (\texttt{3382.T})
    \item Alimentation Couche-Tard (\texttt{ATD.TO})
    \item FEMSA (\texttt{FEMSAUBD.MX})
    \item BP (\texttt{BP.L})
    \item Carrefour (\texttt{CA.PA})
\end{itemize}

\section{Metodología}
\subsection{Análisis técnico}
Se calcularon SMA, EMA, RSI, MACD, Bandas de Bollinger y Estocástico.

\subsection{Rendimientos}
Se estimaron rendimientos simples y logarítmicos, junto con estadísticos descriptivos y pruebas de normalidad.

\subsection{Volatilidad condicional}
Se ajustaron modelos ARCH(1), GARCH(1,1) y EGARCH(1,1), comparados mediante AIC y BIC.

\subsection{CAPM}
Se estimó beta para cada activo frente a un benchmark local y se calculó el rendimiento esperado con CAPM.

\subsection{VaR y CVaR}
Se emplearon tres métodos:
\begin{itemize}
    \item Paramétrico
    \item Histórico
    \item Monte Carlo
\end{itemize}

\subsection{Markowitz}
Se simularon 10{,}000 portafolios para construir el conjunto factible y la frontera eficiente.

\subsection{Señales automáticas}
Se implementaron reglas de compra/venta usando MACD, RSI, Bollinger, cruces de medias y Estocástico.

\subsection{Benchmark y contexto macro}
Se comparó el portafolio contra un benchmark global y se reportaron métricas como Sharpe, Alpha de Jensen, Tracking Error, Information Ratio y máximo drawdown.

\section{Resultados}
En esta sección deben insertarse tablas y figuras del dashboard:
\begin{itemize}
    \item Figura de precios e indicadores.
    \item Tabla descriptiva de rendimientos.
    \item Comparación de modelos GARCH.
    \item Tabla de VaR/CVaR.
    \item Frontera eficiente y pesos óptimos.
    \item Señales por activo.
    \item Desempeño contra benchmark.
\end{itemize}

\section{Conclusiones}
El dashboard permite integrar diferentes enfoques de medición del riesgo en una sola herramienta interactiva. Esto facilita comparar activos, medir exposición sistemática, cuantificar pérdidas extremas y construir portafolios más eficientes.

\section*{Repositorio y ejecución}
\begin{itemize}
    \item Repositorio: completar URL
    \item Ejecución: \texttt{streamlit run app.py}
\end{itemize}

\end{document}
